% =============================================================================
%  Section 1 — Head-to-Head Comparison: MCP v2 vs Raw v3
%  Core section of the evaluation report
% =============================================================================

\section{Head-to-Head Comparison: MCP v2 vs Raw v3}\label{sec:comparison}

\subsection{Setup}

Both runs evaluate the same 22 \texttt{astropy/astropy} tasks from
\texttt{princeton-nlp/SWE-bench\_Verified} using
Claude Opus~4.6 (\texttt{claude-opus-4-6}).
The only controlled variable is the \emph{access method}:

\begin{itemize}
  \item \textbf{MCP v2} --- the model accesses the codebase exclusively
        through ByteBell MCP tools (semantic search, graph traversal,
        file retrieval via knowledge graph).  No local filesystem access.
  \item \textbf{Raw v3} --- the model has direct filesystem access to a
        locally cloned repository at the base commit.
        No MCP tools, no web search, no \texttt{git} CLI.
\end{itemize}

\noindent
Both use the same 8-point rubric
(RC\,3 + Files\,2 + Patch\,3) and are judged by a separate
Claude Opus~4.6 session against the ground-truth patch.

% ---------------------------------------------------------------------------
\subsection{Executive Summary}

\begin{table}[h]
\centering
\caption{Top-line comparison across all 22 tasks.}\label{tab:exec-summary}
\begin{tabular}{llll}
\toprule
\textbf{Metric} & \textbf{MCP v2} & \textbf{Raw v3} & \textbf{Winner} \\
\midrule
Overall score      & \textbf{163/176 (92.6\%)} & 162/176 (92.0\%) & MCP (+1\,pt) \\
Avg cost / task    & \textbf{\$0.52}            & \$0.73            & MCP (1.4$\times$ cheaper) \\
Avg time / task    & 249\,s                     & 251\,s            & Tie \\
Root cause         & 97.0\%                     & 97.0\%            & Tie \\
File ID            & 100\%                      & 100\%             & Tie \\
Patch quality      & \textbf{83.3\%}            & 81.8\%            & MCP \\
Perfect scores     & 13                         & \textbf{14}       & Raw \\
Failures           & \textbf{0}                 & 1                 & MCP \\
\bottomrule
\end{tabular}
\end{table}

\noindent
MCP edges Raw by one point at 29\% lower cost per task.
Both exceed 92\% accuracy.
MCP eliminates all outright failures; Raw produces one more perfect
8/8 score.

% ---------------------------------------------------------------------------
\subsection{Per-Task Head-to-Head}\label{sec:h2h}

Table~\ref{tab:h2h} shows every task.
Bold scores mark disagreements.

\begin{table}[h]
\centering
\small
\caption{Per-task score comparison.  Bold = the two runs disagree.}\label{tab:h2h}
\setlength{\tabcolsep}{3.5pt}
\begin{tabular}{clccl}
\toprule
\textbf{\#} & \textbf{Instance ID} & \textbf{MCP} & \textbf{Raw} & \textbf{Delta} \\
\midrule
 1 & \texttt{12907}  & 8 & 8 & Tie \\
 2 & \texttt{13033}  & 8 & 8 & Tie \\
 3 & \texttt{13236}  & 8 & 8 & Tie \\
 4 & \texttt{13398}  & 8 & 8 & Tie \\
 5 & \texttt{13453}  & \textbf{7} & \textbf{8} & Raw\,+1 \\
 6 & \texttt{13579}  & \textbf{6} & \textbf{8} & Raw\,+2 \\
 7 & \texttt{13977}  & \textbf{8} & \textbf{6} & MCP\,+2 \\
 8 & \texttt{14096}  & 7 & 7 & Tie \\
 9 & \texttt{14182}  & 8 & 8 & Tie \\
10 & \texttt{14309}  & 7 & 7 & Tie \\
11 & \texttt{14365}  & 7 & 7 & Tie \\
12 & \texttt{14369}  & 8 & 8 & Tie \\
13 & \texttt{14508}  & 8 & 8 & Tie \\
14 & \texttt{14539}  & 8 & 8 & Tie \\
15 & \texttt{14598}  & \textbf{5} & \textbf{8} & Raw\,+3 \\
16 & \texttt{14995}  & 8 & 8 & Tie \\
17 & \texttt{7166}   & 8 & 8 & Tie \\
18 & \texttt{7336}   & 8 & 8 & Tie \\
19 & \texttt{7606}   & \textbf{6} & \textbf{3} & MCP\,+3 \\
20 & \texttt{7671}   & 7 & 7 & Tie \\
21 & \texttt{8707}   & \textbf{8} & \textbf{6} & MCP\,+2 \\
22 & \texttt{8872}   & 7 & 7 & Tie \\
\midrule
   & \textbf{Total}  & \textbf{163} & \textbf{162} & \textbf{MCP\,+1} \\
\bottomrule
\end{tabular}
\end{table}

\noindent
16 tasks tie.
6 tasks diverge --- MCP gains +7 points across three tasks,
Raw gains +6 points across three tasks, yielding a net MCP\,+1.

% ---------------------------------------------------------------------------
\subsection{Dimension Breakdown}

\begin{table}[h]
\centering
\caption{Scoring dimensions across all 22 tasks.}\label{tab:dim-cmp}
\begin{tabular}{lrrrrl}
\toprule
\textbf{Dimension} & \textbf{MCP} & \textbf{Raw} & \textbf{MCP\,\%} & \textbf{Raw\,\%} & \textbf{Delta} \\
\midrule
Root cause   & 64/66 & 64/66 & 97.0\% & 97.0\% & Tie \\
Correct files & 44/44 & 44/44 & 100\%  & 100\%  & Tie \\
Correct patch & 55/66 & 54/66 & 83.3\% & 81.8\% & MCP\,+1 \\
\midrule
\textbf{Overall} & \textbf{163/176} & \textbf{162/176} & \textbf{92.6\%} & \textbf{92.0\%} & \textbf{MCP\,+1} \\
\bottomrule
\end{tabular}
\end{table}

\noindent
Root cause and file identification are identical.
The entire 1-point gap comes from patch quality.
MCP wins on tasks 13977, 7606, and 8707 (+4 patch points);
Raw wins on 13453, 13579, and 14598 ($-$3 patch points).
Net: MCP\,+1 patch point.

% ---------------------------------------------------------------------------
\subsection{Resource Usage at a Glance}

\begin{table}[h]
\centering
\caption{Total resource usage across all 22 tasks.}\label{tab:resources}
\begin{tabular}{lrrrrr}
\toprule
& \textbf{Score} & \textbf{Input Tokens} & \textbf{Output Tokens} & \textbf{Time} & \textbf{Cost} \\
\midrule
\textbf{MCP v2} & 163/176 & 1.72\,M & 126\,K & 91\,min & \$11.53 \\
\textbf{Raw v3} & 162/176 & 2.87\,M & 277\,K & 92\,min & \$16.16 \\
\midrule
\textbf{Raw / MCP} & ---   & \textbf{1.7$\times$} & \textbf{2.2$\times$} & 1.0$\times$ & \textbf{1.4$\times$} \\
\bottomrule
\end{tabular}
\\[4pt]
{\footnotesize Input Tokens = cost-weighted: Input + $1.25 \times$ Cache Write + $0.1 \times$ Cache Read.}
\end{table}

\begin{table}[h]
\centering
\caption{Cost of each score point.}\label{tab:efficiency}
\begin{tabular}{lrrrr}
\toprule
& \textbf{Tokens / pt} & \textbf{Output / pt} & \textbf{Time / pt} & \textbf{Cost / pt} \\
\midrule
\textbf{MCP v2} & \textbf{10{,}540}  & \textbf{773} & 33.6\,s & \textbf{\$0.071} \\
\textbf{Raw v3} & 17{,}694            & 1{,}710      & 34.1\,s & \$0.100 \\
\bottomrule
\end{tabular}
\end{table}

\noindent
\textbf{In plain terms:} MCP reads 40\% less code and produces
half the output to achieve the same accuracy.
Both take the same wall-clock time ($\approx$\,91\,min for 22 tasks).
MCP costs \$0.07 per score point; Raw costs \$0.10 --- a 29\% saving.

% ---------------------------------------------------------------------------
\subsection{Side-by-Side Per-Task Comparison}\label{sec:side-by-side}

The following tables place every task's MCP and Raw metrics next to each
other for direct inspection.  Rows where the two runs disagree on score
are shaded.

The input metric used throughout is \textbf{Input Tokens (Weighted)}
--- a single number that abstracts away raw input, cache write, and
cache read tokens into one cost-aligned figure.
Each token class is weighted by its Anthropic pricing tier:
full-price input at $1\times$, cache writes at $1.25\times$, and
cache reads at $0.1\times$.
This means 100K weighted input tokens $\approx$ the cost of reading
100K fresh input tokens, regardless of how many came from cache.

% --- Table A: Score, Weighted Input, Output, Time, Cost ---
\begin{table}[!htbp]
\centering
\small
\caption{Per-task comparison --- Score, Input Tokens (Weighted), Output Tokens, Time, Cost.
         Shaded rows = score disagreement.}\label{tab:sbs-main}
\setlength{\tabcolsep}{3.5pt}
\begin{tabular}{l cc rr rr rr rr}
\toprule
& \multicolumn{2}{c}{\textbf{Score}}
& \multicolumn{2}{c}{\textbf{Input (Weighted)}}
& \multicolumn{2}{c}{\textbf{Output}}
& \multicolumn{2}{c}{\textbf{Time (s)}}
& \multicolumn{2}{c}{\textbf{Cost (\$)}} \\
\cmidrule(lr){2-3}\cmidrule(lr){4-5}\cmidrule(lr){6-7}\cmidrule(lr){8-9}\cmidrule(lr){10-11}
\textbf{Task}
& \textbf{MCP} & \textbf{Raw}
& \textbf{MCP} & \textbf{Raw}
& \textbf{MCP} & \textbf{Raw}
& \textbf{MCP} & \textbf{Raw}
& \textbf{MCP} & \textbf{Raw} \\
\midrule
\texttt{12907}  & 8 & 8 & 53K & 141K & 1.5K & 7.8K & 43 & 259 & 0.30 & 0.40 \\
\texttt{13033}  & 8 & 8 & 38K & 79K & 2.5K & 11.6K & 66 & 191 & 0.24 & 0.67 \\
\texttt{13236}  & 8 & 8 & 98K & 233K & 4.5K & 26.2K & 131 & 422 & 0.60 & 1.40 \\
\texttt{13398}  & 8 & 8 & 121K & 317K & 4.7K & 35.0K & 281 & 612 & 0.70 & 1.77 \\
\rowcolor{gray!15}
\texttt{13453}  & \textbf{7} & \textbf{8} & 62K & 106K & 2.4K & 5.2K & 71 & 123 & 0.34 & 0.64 \\
\rowcolor{gray!15}
\texttt{13579}  & \textbf{6} & \textbf{8} & 74K & 84K & 5.7K & 13.3K & 472 & 223 & 0.50 & 0.74 \\
\rowcolor{gray!15}
\texttt{13977}  & \textbf{8} & \textbf{6} & 74K & 114K & 4.7K & 9.3K & 392 & 212 & 0.48 & 0.79 \\
\texttt{14096}  & 7 & 7 & 65K & 134K & 4.8K & 9.0K & 134 & 175 & 0.44 & 0.62 \\
\texttt{14182}  & 8 & 8 & 80K & 172K & 3.2K & 10.5K & 82 & 232 & 0.48 & 0.57 \\
\texttt{14309}  & 7 & 7 & 49K & 24K & 1.7K & 1.7K & 52 & 171 & 0.28 & 0.15 \\
\texttt{14365}  & 7 & 7 & 76K & 64K & 2.9K & 7.3K & 96 & 141 & 0.45 & 0.50 \\
\texttt{14369}  & 8 & 8 & 166K & 144K & 35.9K & 31.4K & 845 & 475 & 1.72 & 1.49 \\
\texttt{14508}  & 8 & 8 & 56K & 107K & 2.2K & 5.4K & 413 & 100 & 0.33 & 0.39 \\
\texttt{14539}  & 8 & 8 & 57K & 73K & 2.2K & 6.2K & 416 & 123 & 0.33 & 0.51 \\
\rowcolor{gray!15}
\texttt{14598}  & \textbf{5} & \textbf{8} & 202K & 192K & 18.8K & 32.6K & 547 & 600 & 1.47 & 1.77 \\
\texttt{14995}  & 8 & 8 & 57K & 132K & 1.5K & 6.8K & 55 & 153 & 0.31 & 0.51 \\
\texttt{7166}   & 8 & 8 & 84K & 122K & 7.5K & 7.4K & 145 & 151 & 0.60 & 0.39 \\
\texttt{7336}   & 8 & 8 & 58K & 47K & 3.0K & 3.1K & 83 & 291 & 0.36 & 0.31 \\
\rowcolor{gray!15}
\texttt{7606}   & \textbf{6} & \textbf{3} & 59K & 96K & 2.5K & 6.9K & 94 & 181 & 0.35 & 0.39 \\
\texttt{7671}   & 7 & 7 & 65K & 75K & 7.0K & 12.6K & 785 & 202 & 0.49 & 0.68 \\
\rowcolor{gray!15}
\texttt{8707}   & \textbf{8} & \textbf{6} & 63K & 320K & 3.2K & 20.7K & 84 & 372 & 0.39 & 1.14 \\
\texttt{8872}   & 7 & 7 & 61K & 91K & 3.6K & 7.0K & 187 & 119 & 0.39 & 0.35 \\
\midrule
\textbf{Total}  & \textbf{163} & \textbf{162}
                & \textbf{1.72M} & \textbf{2.87M}
                & \textbf{126K} & \textbf{277K}
                & \textbf{5{,}474} & \textbf{5{,}528}
                & \textbf{11.53} & \textbf{16.16} \\
\textbf{Avg}    & \textbf{7.4} & \textbf{7.4}
                & \textbf{78K} & \textbf{130K}
                & \textbf{5.7K} & \textbf{12.6K}
                & \textbf{249} & \textbf{251}
                & \textbf{0.52} & \textbf{0.73} \\
\bottomrule
\end{tabular}
\\[4pt]
{\footnotesize
M = MCP, R = Raw.\quad
Input (Weighted) = Input + $1.25 \times$ CacheWrite + $0.1 \times$ CacheRead.\quad
``K'' = thousands, ``M'' = millions of tokens.}
\end{table}

\noindent
\textbf{Key patterns in the per-task data:}

\begin{itemize}[nosep, leftmargin=*]
\item \textbf{MCP reads fewer input tokens on 18 of 22 tasks.}
      The 4 exceptions where Raw is leaner:
      \texttt{14309} (24K vs 49K --- Raw's lightest task),
      \texttt{14365} (64K vs 76K),
      \texttt{14369} (144K vs 166K),
      and \texttt{14598} (192K vs 202K --- the sole efficiency inversion,
      see \S\ref{sec:anomalies}).
      All 4 tie on score --- Raw's occasional input advantage does not
      translate to accuracy gains.

\item \textbf{Output tokens are 2.2$\times$ higher in Raw overall.}
      Raw produces more output on 20/22 tasks.
      The two exceptions (\texttt{14369}: 31.4K Raw vs 35.9K MCP, and
      \texttt{7166}: 7.4K vs 7.5K) are near parity.

\item \textbf{Wall-clock time shows no consistent winner.}
      MCP is faster on 11 tasks, Raw is faster on 11 tasks.
      Totals are near-identical (5{,}474\,s vs 5{,}528\,s).
      Tasks \texttt{7671} (MCP 785\,s vs Raw 202\,s) and
      \texttt{14508} (MCP 413\,s vs Raw 100\,s) are MCP's slowest
      outliers; \texttt{13398} (Raw 612\,s vs MCP 281\,s) and
      \texttt{13236} (Raw 422\,s vs MCP 131\,s) are Raw's.

\item \textbf{Cost tracks input tokens, not time.}
      MCP is cheaper on 16/22 tasks.  The 6 exceptions
      (\texttt{14309}, \texttt{14365}, \texttt{14369}, \texttt{7166},
      \texttt{7336}, \texttt{8872}) are all tasks where Raw's
      Haiku-heavy exploration is very cheap per token.

\item \textbf{Largest divergence: task 8707.}
      Raw consumed 320K weighted input tokens (5.1$\times$ MCP's 63K)
      and 372\,s (4.4$\times$ MCP's 84\,s), yet scored 6/8 versus
      MCP's 8/8.  This is the benchmark's most extreme case of
      exploration without synthesis.
\end{itemize}

% --- Table: How each method uses its two models ---
\begin{table}[!htbp]
\centering
\caption{How each method splits work between Opus (reasoning) and Haiku (lightweight).}\label{tab:model-cmp}
\begin{tabular}{l cc}
\toprule
& \textbf{MCP v2} & \textbf{Raw v3} \\
\midrule
\textbf{Opus cost}          & \$11.48 (99.5\%)   & \$14.89 (92.1\%) \\
\textbf{Haiku cost}         & \$0.05\ \ (0.5\%)  & \$1.27\ \ (7.9\%) \\
\textbf{Opus API calls}     & 266                & 350 \\
\textbf{Haiku API calls}    & 22 (1 per task)    & 178 (avg 8 per task) \\
\textbf{Haiku role}         & Route \& hand off  & Explore codebase \\
\bottomrule
\end{tabular}
\end{table}

\noindent
\textbf{In plain terms:} Claude Code uses two models internally ---
Opus (expensive, does the thinking) and Haiku (cheap, does lightweight
tasks).  The two methods use them very differently:

\begin{itemize}[nosep, leftmargin=*]
\item \textbf{Under MCP}, Haiku makes exactly 1 call per task just to
      route the request, then Opus does everything through the
      knowledge graph.  99.5\% of the cost is Opus.

\item \textbf{Under Raw}, Haiku is an active explorer --- it makes
      up to 39 calls per task to search files, list directories, and
      read code.  This multi-agent exploration is where Raw's extra
      token volume (and 8\% of its cost) comes from.

\item \textbf{Opus itself is 23\% cheaper under MCP} (\$11.48 vs
      \$14.89) because the knowledge graph pre-filters relevant code,
      so Opus needs fewer reasoning turns.
\end{itemize}

% ---------------------------------------------------------------------------
\subsection{Why Token Ratios $\neq$ Cost Ratios}
\label{sec:token-cost-gap}

The cache-read ratio (1.9$\times$) exceeds the cost ratio
(1.4$\times$) because Raw's extra tokens are
disproportionately cheap Haiku sub-agent tokens.

\begin{table}[h]
\centering
\caption{Token ratio vs.\ cost ratio --- the model-mix discount.}\label{tab:ratio-gap}
\begin{tabular}{lrl}
\toprule
\textbf{Metric} & \textbf{Ratio (Raw/MCP)} & \textbf{Why it differs from cost} \\
\midrule
Cache Read tokens          & 1.9$\times$ & Cache reads priced at $0.1\times$ input rate \\
Effective Input (unwt.)    & 1.8$\times$ & Treats all input tokens as equal cost \\
\textbf{Eff Weighted Input}& \textbf{1.7$\times$} & \textbf{Weights by tier --- closest to cost} \\
Actual cost                & 1.4$\times$ & Output costs + haiku/opus model mix \\
\bottomrule
\end{tabular}
\end{table}

\noindent
Raw v3 uses multi-agent exploration where Haiku sub-agents perform
file search, directory listing, and code reading, generating the bulk
of Raw's cache reads at Haiku rates (\$0.10/MTok --- 5$\times$ cheaper
than Opus).  MCP v2 routes nearly all work through Opus directly.
The Effective Weighted Input metric (1.7$\times$) tracks closest to
actual cost because it encodes Anthropic's pricing tiers into the
formula.  The residual gap to 1.4$\times$ is the Haiku model-mix
discount.

% ---------------------------------------------------------------------------
\subsection{The 6 Disagreement Tasks --- Detailed Analysis}
\label{sec:disagreements}

\subsubsection{Raw wins: \texttt{14598} (+3), \texttt{13579} (+2),
               \texttt{13453} (+1)}

\paragraph{Task 15 (\texttt{14598}) --- Raw +3.}
FITS CONTINUE cards lose quotes from double un-escaping.
The ground truth fixes the \emph{parser side}: an end-of-string anchor
on \texttt{\_strg\_comment\_RE} and removal of
\texttt{.replace("''","'")} in \texttt{\_split()}.
MCP proposed a \emph{writer-side} fix
(\texttt{\_format\_long\_image}) --- correct symptom identification but
wrong causal direction.
Raw traced the double-unescaping flow through direct file reads and
matched the GT exactly.

\emph{Data note:} This is MCP's costliest task (\$1.47) and
second-highest EW (202K), yet produced its lowest score (5/8).
High exploration cost did not translate to accuracy.
Raw spent comparably (\$1.77, EW 192K) but achieved 8/8.
\textbf{This is the only task where MCP's EW exceeds Raw's} (202K vs
192K), making it an anomalous efficiency inversion.

\paragraph{Task 6 (\texttt{13579}) --- Raw +2.}
\texttt{SlicedLowLevelWCS.world\_to\_pixel\_values} uses a hardcoded
\texttt{1.0} for dropped dimensions.
The GT is a single-line addition.
MCP proposed a multi-step two-pass iterative approach ---
over-engineered relative to the GT's simplicity.
Raw found the minimal fix.

\emph{Data note:} MCP spent 472\,s (2.1$\times$ Raw's 223\,s)
on this task, suggesting the two-pass approach emerged from extended
exploration rather than a quick diagnosis.  Despite the extra time,
the more complex solution scored lower.

\paragraph{Task 5 (\texttt{13453}) --- Raw +1.}
HTML writer ignores the \texttt{formats} argument.
Both identify the two required lines
(\texttt{self.data.cols = cols} + \texttt{\_set\_col\_formats()}).
MCP's insertion point differs slightly from the GT, resulting in a
functionally equivalent but not exact match.
Raw places both lines precisely.

\emph{Data note:} The smallest disagreement in the benchmark.
MCP used 1.9$\times$ fewer EW tokens (62K vs 106K) yet lost by 1
point, showing that token efficiency does not guarantee patch
precision.

\subsubsection{MCP wins: \texttt{7606} (+3), \texttt{13977} (+2),
               \texttt{8707} (+2)}

\paragraph{Task 19 (\texttt{7606}) --- MCP +3.}
\texttt{Unit.\_\_eq\_\_} should return \texttt{NotImplemented} instead
of \texttt{False}.
The GT requires changes in two locations:
\texttt{UnitBase.\_\_eq\_\_} and \texttt{UnrecognizedUnit.\_\_eq\_\_}.
MCP fixed one of the two (\texttt{UnrecognizedUnit}), scoring 6/8.
Raw returned only a filename with no root cause and no patch,
scoring 3/8 --- its worst result across all 22 tasks.

\emph{Data note:} Raw's failure is anomalous --- task 7606 is rated
\emph{15\,min--1\,h} difficulty, yet Raw spent only \$0.39 and 181\,s.
This is below its average cost (\$0.73) and time (251\,s),
suggesting premature termination rather than inability.
MCP spent even less (\$0.35, 94\,s) and still produced a partial fix.

\paragraph{Task 7 (\texttt{13977}) --- MCP +2.}
\texttt{Quantity.\_\_array\_ufunc\_\_} should return
\texttt{NotImplemented} for duck types.
The GT wraps the \emph{entire} method body in
\texttt{try/except} and checks other operands'
\texttt{\_\_array\_ufunc\_\_} before returning
\texttt{NotImplemented}.
MCP matched this approach exactly.
Raw wrapped only the input conversion loop (3 lines) --- too narrow.

\emph{Data note:} MCP's EW (74K) is 35\% lower than Raw's (114K) on
this task, yet MCP achieved the better score.  The knowledge graph's
semantic retrieval of the full method context appears to have guided
the model toward the complete fix, while Raw's line-by-line
exploration led to a partial one.

\paragraph{Task 21 (\texttt{8707}) --- MCP +2.}
\texttt{Card.fromstring} / \texttt{Header.fromstring} should accept
\texttt{bytes}.
MCP produced complete code: \texttt{Card} decodes latin1,
\texttt{Header} handles bytes-aware CONTINUE/END/separator.
Raw described the fix in prose but generated no actual diff.

\emph{Data note:} Raw spent 4.4$\times$ the time (372\,s vs 84\,s)
and 5.1$\times$ the EW tokens (320K vs 63K), yet scored 2 points
lower.  This is the largest efficiency gap per score point in the
benchmark: Raw burned 320K EW tokens for a 6/8 answer while MCP
achieved 8/8 with 63K.  Raw's multi-agent exploration generated
extensive context (1.94M cache reads) but failed to crystallise it
into code.

% ---------------------------------------------------------------------------
\subsection{Shared Failure Patterns}\label{sec:shared}

Three tasks tie at 7/8 with the \emph{same} deviation from the
ground truth (Table~\ref{tab:shared-fail}).  This is notable because
the deviations are identical despite completely different access
methods, suggesting the alternative implementations are a property of
the model's reasoning rather than the information source.

\begin{table}[h]
\centering
\caption{Tasks where both runs make the same mistake.}\label{tab:shared-fail}
\begin{tabular}{lp{9cm}}
\toprule
\textbf{Task} & \textbf{Shared deviation from GT} \\
\midrule
\texttt{14365} & Both add \texttt{re.IGNORECASE} but miss \texttt{v.upper() == "NO"} for data-value parsing. \\
\texttt{14309} & Both use a \texttt{bool(args)} guard instead of GT's structural \texttt{return} flow fix. \\
\texttt{14096} & Both walk the MRO to find the property descriptor instead of GT's simpler \texttt{\_\_getattribute\_\_} re-invoke. \\
\bottomrule
\end{tabular}
\end{table}

\noindent
These shared misses indicate that when the GT's approach is
non-obvious, the model converges on the \emph{same} alternative
implementation regardless of whether it reads code via MCP or
directly from disk.  The information bottleneck is in the model's
reasoning, not the retrieval path.

% ---------------------------------------------------------------------------
\subsection{Complementary Failure Modes}\label{sec:complementary}

The two approaches fail in non-overlapping ways
(Table~\ref{tab:fail-modes}).

\begin{table}[h]
\centering
\caption{Failure modes by access method.}\label{tab:fail-modes}
\begin{tabular}{lll}
\toprule
\textbf{Failure type} & \textbf{MCP v2} & \textbf{Raw v3} \\
\midrule
Wrong causal path (writer vs parser) & 14598 (5/8)  & --- \\
Over-engineering a simple fix        & 13579 (6/8)  & --- \\
No code output (description only)    & ---           & 7606 (3/8), 8707 (6/8) \\
Too-narrow patch scope               & ---           & 13977 (6/8) \\
\bottomrule
\end{tabular}
\end{table}

\noindent
MCP fails when the knowledge graph leads the model down a plausible
but incorrect causal thread, or when the model over-engineers from
broad context.
Raw fails when extensive file-level exploration does not
crystallise into code output --- the model ``understands'' the fix but
stops short of writing it.

A hypothetical \textbf{best-of-both ensemble} (taking the higher score
per task) would achieve \textbf{170/176 (96.6\%)}, recovering 7--8
points that each approach loses individually.

% ---------------------------------------------------------------------------
\subsection{Anomalies \& Data Notes}\label{sec:anomalies}

The following irregularities were observed in the data.  All are
derived from the per-task metrics in the source reports.

\begin{enumerate}

\item \textbf{Task 14598 is an efficiency inversion.}
      MCP's EW (202K) \emph{exceeds} Raw's (192K) --- the only task
      where this occurs.  On every other task, MCP uses fewer EW tokens.
      Despite spending more tokens, MCP scored 3 points lower (5 vs 8).
      This suggests the knowledge graph surfaced a large volume of
      FITS card-handling code but biased the model toward the writer
      path rather than the parser path.

\item \textbf{Raw's 7606 failure is anomalous for its difficulty.}
      Task 7606 is rated $<$15\,min--1\,h and requires a 2-line change
      (\texttt{return NotImplemented} in two \texttt{\_\_eq\_\_}
      methods).  Raw spent below-average resources (\$0.39, 181\,s)
      and produced no patch.  This appears to be premature
      termination --- the model decided it had answered (returning
      only a filename) when it had not.  MCP, also spending below
      average (\$0.35, 94\,s), at least produced a half-fix.

\item \textbf{Raw's 8707 is the largest EW waste.}
      Raw consumed 320K EW tokens (5.1$\times$ MCP's 63K) and 372\,s
      (4.4$\times$ MCP's 84\,s) on task 8707, yet scored 6/8 versus
      MCP's 8/8.  The 1.94M cache-read tokens indicate deep
      multi-agent exploration of \texttt{card.py} and
      \texttt{header.py} that did not produce a diff.
      This is the benchmark's most extreme case of exploration
      without synthesis.

\item \textbf{Cost and score are weakly correlated.}
      Across both runs, higher cost does not reliably predict higher
      score.  MCP's cheapest task (\texttt{13033}, \$0.24) scores 8/8;
      its costliest (\texttt{14598}, \$1.47) scores 5/8.
      Raw's cheapest (\texttt{14309}, \$0.15) scores 7/8; its
      costliest (\texttt{13398}, \$1.77) scores 8/8 but so does
      \texttt{14598} at the same price.
      The Pearson correlation between cost and score is near zero for
      both runs.

\item \textbf{Wall-clock time is indistinguishable between methods.}
      Total time is 5{,}474\,s (MCP) vs 5{,}528\,s (Raw) --- a 1\%
      difference.  Despite MCP using 1.7$\times$ fewer EW tokens,
      the wall-clock savings are absorbed by MCP tool-call round-trip
      latency.  Token efficiency translates to dollar savings but not
      to time savings at this scale.

\item \textbf{Haiku's role differs sharply between methods.}
      In MCP, Haiku makes exactly 1 API request per task (22/22) with
      zero cache tokens --- a pure routing stub, accounting for
      \$0.052 total (0.5\% of cost).
      In Raw, Haiku is an active exploration agent: up to 39 requests
      on task 8707, with 4.3M total cache-read tokens, accounting for
      \$1.27 (7.9\% of cost).
      This means MCP's token budget is almost entirely Opus-priced,
      while Raw's includes a substantial volume of cheap Haiku tokens
      --- explaining why MCP's 1.7$\times$ EW advantage compresses to
      only 1.4$\times$ cost advantage.

\item \textbf{Three shared 7/8 misses are model-intrinsic.}
      Tasks 14096, 14309, and 14365 produce the same alternative
      implementation under both access methods.  Since the only
      controlled variable is the retrieval path, these deviations
      originate in the model's reasoning preferences, not in
      information availability.  This puts a floor on accuracy
      improvements achievable by better retrieval alone: these 3
      points ($\approx$1.7\% of 176) require changes to the model
      itself or to the prompt, not to the knowledge graph.

\item \textbf{The ensemble ceiling is 96.6\%.}
      Taking the higher score per task yields 170/176.  The
      remaining 6 lost points come from tasks where \emph{both}
      approaches score $<$8: tasks 14096 (7), 14309 (7), 14365 (7),
      7671 (7), and 8872 (7).  All are the shared 7/8 pattern ---
      confirming that retrieval improvements cannot close this gap.

\end{enumerate}

% ---------------------------------------------------------------------------
\subsection{Verdict}\label{sec:verdict}

\begin{table}[h]
\centering
\caption{Final verdict across all dimensions.}\label{tab:verdict}
\begin{tabular}{lll}
\toprule
\textbf{Dimension} & \textbf{Winner} & \textbf{Detail} \\
\midrule
Overall accuracy       & MCP   & 163/176 (92.6\%) vs 162/176 (92.0\%) \\
Root cause             & Tie   & Both 97.0\% \\
File identification    & Tie   & Both 100\% \\
Patch correctness      & MCP   & 83.3\% vs 81.8\% \\
Exact 8/8 count        & Raw   & 14 vs 13 \\
Catastrophic failures  & MCP   & 0 vs 1 \\
Cost                   & MCP   & \$11.53 vs \$16.16 (1.4$\times$) \\
Time                   & Tie   & 5{,}474\,s vs 5{,}528\,s \\
EW per point           & MCP   & 10{,}540 vs 17{,}694 (1.7$\times$) \\
\bottomrule
\end{tabular}
\end{table}

\noindent
MCP edges Raw by 1 point at 29\% lower cost per task.
Both exceed 92\% accuracy with identical root-cause and
file-identification performance.
The 1-point overall gap comes entirely from patch quality.
Their failure modes are complementary --- MCP struggles with causal
misattribution and over-engineering; Raw struggles with answer
finalisation.
An ensemble taking the best of each would reach 96.6\%.

\paragraph{Practical recommendation.}
\begin{itemize}
  \item \textbf{Cost-sensitive production:} MCP v2 --- same or better
        accuracy at 1.4$\times$ fewer dollars.
  \item \textbf{Maximum accuracy:} Ensemble best-of-both --- projected
        96.6\%.
  \item \textbf{Maximum perfect scores:} Raw v3 --- 14 exact
        matches vs 13, at the cost of one 3/8 failure.
\end{itemize}
